% Pursuing Stacks (cotes 134-2 to 134-8): archive cote 134-5, site page 129 (PDF page 130), typescript page 316.
% typescript. Draft compared: ps5.tex:1171-1256.
% AI-transcribed by Claude Opus 5.5 (claude-opus-5-5, adaptive thinking, effort high), Message Batches API, batch ps3 (msgbatch_01XdFmuVVBAcgc7EFxdtCqsj), 2026-11-12; not yet checked by eye.
\scanpage{134-5/129}
\mspage{316}

It should'nt be hard,  I feel,  to pin down exactly what is needed for
getting such a formalism.  The intuitive idea behind it (suggested by
the example of \xout{the} standard complexes and multicomplexes)  is that
\uline{such a formalism should be associated to cellular decompositions of}
$n$-\uline{cells} \uline{for variable} $n$,  such that the interior of each $n$-cell should
be an open cell of the subdivision.  T\unsure{\hand{h}}ere may of course be several
$n$-cells\xout{\unsure{s}} for the same $n$,  which are not combinatorially isomorphic.
When trying to express this idea in a precise way,  we are led to assu-%
me,  as an extra structure on the would-be test category $A$,  a functor
\[ (1) \qquad i:\ A \hand{\longrightarrow} (\mathrm{Ord}) \]
of $A$ into the category of ordered sets,  such that for any $a$ in $A$,
$i(a)$ be a \uline{finite} ordered set,  whose geometric realization  (cf. section
22)  is an $n$-cell for suitable  $n \overset{\hand{def}}{=} \mathrm{dim}(a)$.  We assume moreover that
$i(a)$ has a largest element  $e(a)$,  and that the geometric realization
of  $i(a) \hand{\setminus\{} e(a) \hand{\}}$ is the bounding  $(n-1)$-sphere of the $n$-cell  $\hand{|} i(a) \hand{|}$ :
\[ (2) \qquad \hand{|} i(a)^{*} \hand{|}\ \deleted{\unsure{homemomorphic|homeomorphic} to}\ \hand{\simeq}\ S^{n-1} \hand{,}\ \hand{\text{where}}\ i(a)^{*} \overset{\hand{def}}{=} i(a) \hand{\setminus\{e(a)\}} \hand{,} \]
which will imply the precedent condition namely
\[ (3) \qquad \hand{|} i(a) \hand{|} \hand{\simeq} B^{n} \quad , \]
as $B^{n}$ can be identified\xout{\unsure{e}} to the cone over $S^{n-1}$.  As another condition,
we need that \xout{\unsure{for any}}

(4) \quad For any $a$ in $A$ and  $x \hand{\in} i(a)$,  there exists $b$ in $A$ and an iso-%
morphism
\[ i(b) \hand{\xrightarrow{\sim}} i(a)_{/x} \overset{\mathrm{def}}{=} \hand{\{}\, y \hand{\in} i(a) \hand{|}\ y \hand{\leq} x \,\hand{\}} \ \hand{\hookrightarrow i(a)}, \]
induced by a map
\[ b \hand{\longrightarrow} a \]
in $A$.

It is enough to make this assumption for $x$ of codimension 1 in $i(a)$,
which will imply thatnit is true for any $x$.  It seems reasonable on the
other hand to make the assumption that for a given $x$ ,  the object  $b$
in (4),  viewed as an object of  $A_{/a}$ \hand{,}  is determined up to a unique
isomorphism ,  we may call it $a_{x}$ ,  and \inserted{$\partial_x$} the canonical map of $b$ into $a$
\margin{\textcircled{$a_x$}}
\[ (5) \qquad \hand{\partial}_{x} :\ a_{x} \hand{\longrightarrow} a \qquad ( x \hand{\in} i(a)\hand{,} \text{ of codim.1 in } i(a)\xout{\illegible{1}}\hand{,} \]
\[ \text{i.e. } \mathrm{dim}(x)=\mathrm{dim}(a)-1 \ ) \quad . \]
As an extra structure,  we need for any $a$ in $A$
\[ (6) \qquad \hand{\omega}_{a} \ ,\ \text{an \uline{orientation} of the $n$-cell } i(a) \ \ (n=\mathrm{dim}(a)). \]
This allows us,  for any $x$ as in (5),  to define a signature
\[ (7) \qquad \hand{\varepsilon}(x) \ \text{ or } \ \hand{\varepsilon}_{a}(x) \hand{\in} \hand{\{} +1,-1 \hand{\}} \ , \]

\margin{129}
